%%%%%%%%%%%%%%%%%%%%%%%%%%%%%%%%%%%%%%%%%
% Simple Sectioned Essay Template
% LaTeX Template
%
% This template has been downloaded from:
% http://www.latextemplates.com
%
% Note:
% The \lipsum[#] commands throughout this template generate dummy text
% to fill the template out. These commands should all be removed when 
% writing essay content.
%
%%%%%%%%%%%%%%%%%%%%%%%%%%%%%%%%%%%%%%%%%

%----------------------------------------------------------------------------------------
%	PACKAGES AND OTHER DOCUMENT CONFIGURATIONS
%----------------------------------------------------------------------------------------

\documentclass[12pt]{article} % Default font size is 12pt, it can be changed here
\usepackage[T1]{fontenc}      
\usepackage[swedish]{babel}
\usepackage[utf8]{inputenc}
\usepackage{listings}
\usepackage{geometry} % Required to change the page size to A4
\usepackage{hyperref}
\geometry{a4paper} % Set the page size to be A4 as opposed to the default US Letter

\usepackage{graphicx} % Required for including pictures

\usepackage{float} % Allows putting an [H] in \begin{figure} to specify the exact location of the figure
\usepackage{wrapfig} % Allows in-line images such as the example fish picture

\usepackage{lipsum} % Used for inserting dummy 'Lorem ipsum' text into the template

\linespread{1.2} % Line spacing

%\setlength\parindent{0pt} % Uncomment to remove all indentation from paragraphs

\graphicspath{{Pictures/}} % Specifies the directory where pictures are stored

\begin{document}

%----------------------------------------------------------------------------------------
%	TITLE PAGE
%----------------------------------------------------------------------------------------

\begin{titlepage}

\newcommand{\HRule}{\rule{\linewidth}{0.5mm}} % Defines a new command for the horizontal lines, change thickness here

\center % Center everything on the page

\textsc{\LARGE Uppsala Universitet}\\[1.5cm] % Name of your university/college
\textsc{\Large Civilingenjörsprogrammet i Informationsteknologi}\\[0.5cm] % Major heading such as course name
\textsc{\large Imperativ och Objectorienterad Progammeringingsmetodik}\\[0.5cm] % Minor heading such as course title

\HRule \\[0.4cm]
{ \huge \bfseries Cedet och ECB för Emacs}\\[0.4cm] % Title of your document
\HRule \\[1.5cm]

\begin{minipage}{0.4\textwidth}
\begin{flushleft} \large
\emph{Author:}\\
Max \textsc{Andersson} % Your name
\end{flushleft}
\end{minipage}
~
\begin{minipage}{0.4\textwidth}
\begin{flushright} \large
\emph{Supervisor:} \\
Tobias \textsc{Wrigstad} % Supervisor's Name
\end{flushright}
\end{minipage}\\[4cm]

{\large \today}\\[3cm] % Date, change the \today to a set date if you want to be precise

%\includegraphics{Logo}\\[1cm] % Include a department/university logo - this will require the graphicx package

\vfill % Fill the rest of the page with whitespace

\end{titlepage}

%----------------------------------------------------------------------------------------
%	TABLE OF CONTENTS
%----------------------------------------------------------------------------------------

\tableofcontents % Include a table of contents

\newpage % Begins the essay on a new page instead of on the same page as the table of contents 

%----------------------------------------------------------------------------------------
%	INTRODUCTION
%----------------------------------------------------------------------------------------

\section{Introduction} % Major section

CEDET (Collection of Emacs Development Enviroment Tools) och ECB (Emacs Code Browser) är verktyg utvecklade för berika Emacs som utvecklingsmiljö. % \cite{Figueredo:2009dg}.

%------------------------------------------------

\subsection{Vad är CEDET?} % Sub-section

CEDET är en samling verktyg. Det innehåller olika verktyg som gör en programmerares liv i Emacs lättare. 
Emacs är en textediterare, men med tillägg som CEDET förvandlars det till en utvecklingsmiljö.  I denna utvecklingsmiljö kan vi bla. strukturera, navigera, generera och automatisera vår kod. Cedet har även en inbyggd PMS (Projekt Management System), som gör det möjligt för oss att inte bara strukturera vår kod, utan hela vårt projekt.

%------------------------------------------------

\subsection{Vad är ECB?} % Sub-section

ECB eller Emacs Code Browser är precis vad namnet föreslår, en navigator för att navigera din kod. Det som gör ECB extra bra är det att det finns möjlighet att överblicka strukturen i kod filer, arbeta strukturerat med mappar, ha tillgång till de resurser ditt projekt behöver och ger dig en visuell överblick av det samtidigt som du kodar. Dessutom drar det nytta av Cedet paketets inbyggda delar, som gör ECB's funktionalitet möjligt. 

ECB använder sig av trädstrukturer för att representera filerna i din mapphierarki, och ger dig snabb åtkomst till dina projektfiler och rader i din dokumentstruktur.

%------------------------------------------------

%\begin{figure}[H] % Example image
%\center{\includegraphics[width=0.5\linewidth]{placeholder}}
%\caption{Example image.}
%\label{fig:speciation}
%\end{figure}

%------------------------------------------------


%----------------------------------------------------------------------------------------
%	MAJOR SECTION 1
%----------------------------------------------------------------------------------------
\newpage
\section{CEDET (Collection of Emacs Development Eviroment Tools)} % Major section

Cedet har många olika verktyg som möjligör och förenklar en arbetsprocess utförd i emacs. Det vi ska ta upp i detta dokument är :
\begin{itemize}
\item Installation och enkel konfigurering av CEDET
\item Stöd för projekthantering och Självgenerande Makefiler
\end{itemize}
%------------------------------------------------


\subsection{Installation och Konfiguration av CEDET} % Sub-section
I de nyare versionerna av Emacs finns redan Cedet installerat och har du det installerat du kan hoppa förbi sektionen om installationen,, om du är osäker så kontrollera om det finns i din version genom att tycka M-x cedet-version RET. 
Installationen för CEDET ser lite olika ut beroende på vilket operativsystem du sitter på, men är i båda fallen relativt simpel.
Först måste du ladda ner en utgåva av CEDET som du kan göra på
{http://sourceforge.net/projects/cedet/} och extrahera paketet någonstans på din dator (rekomenderat är att lägga den i en mapp som tex.  ~/.emacs.d/packages/ ).

%------------------------------------------------
\subsubsection{Installation på Unix baserade system (Mac OS X, Linux)}
Efter att ha laddat ner mjukvaran till din dator så öppna upp terminalen och navigera dig till mappen som du nyss extraherade och kör kommandot:
\begin{lstlisting}
make EMACS=emacs
\end{lstlisting}

%------------------------------------------------

\subsubsection{Installation på Windows baserade system}
På ett Windows baserat system fungerar inte make kommandot, men just detsamma så öppna upp Komtandotolken vilket man smidigt kan göra genom att klicka WIN + R sedan skriva `cmd` följt av ENTER. Därefter navigera till den extraherade mappen och köra kommandot:
\begin{lstlisting}
emacs -Q -l cedet-build.el -f cedet-build 
\end{lstlisting}


%\begin{wrapfigure}{l}{0.4\textwidth} % Inline image example
% \begin{center}
%  \includegraphics[width=0.38\textwidth]{fish}
%  \end{center}
%  \caption{Fish}
%\end{wrapfigure}


%------------------------------------------------

\subsubsection{Konfigurations av Emacs} % Sub-sub-section
När installationen av de utökade verktygsprogrammen är installerade så behöver vi berätta för Emacs att det kan vara bra att starta tilläggen när vi startar Emacs. Detta gör vi genom att lägga till ett par rader kod i våran initieringsfil. Beroende på din version och platform så kan den filen finnas olika ställen men det vanligaste är att du har en .emacs fil i din hemkatalog (`~/.emacs`), men den kanske också vara (`.emacs.d/init.el`). Finns inte filen brukar det gå bra att bara skapa dem.  Tänk på att i Unix system så är filer och mappar som börjar men `.` dolda i filsystemet så det underlättar att öppna filen i terminalen/komandotolken eller med en absolut/relativ sökväg i emacs. 
När du hittat och öppnat filen så klista in eller skriv följande rader i filen
\begin{lstlisting}
;; CEDET Config
  (load-file "~/.emacs.d/packages/cedet-1.1/common/cedet.el")
  (semantic-load-enable-code-helpers) 
  (global-srecode-minor-mode 1)       
  (global-ede-mode t)
\end{lstlisting}

Det finns självklart många fler konfigurationer som går att göra, men detta räcker för de ändamål vi har. Efter att du är klar med att sätta in raderna i filen, så spara filen och starta om emacs eller tryck och skriv  ( M-x eval-buffer ) 
\newpage
\subsection{Stöd Projekthantering och Självgenerande Makefiler}
Cedet är ett mångsidigt program som även fyller funktionen att göra dina Makefiler åt dig. EDE (Emacs Development Enviroment) är det delprogrammet som ingår i Cedet som hanterar projekten. Det första du behöver göra är att ställa dig i en buffer med ett C/C++ dokument där du vill skapa projektet och sedan använda M-x ede-new RET för att skapa ett nytt projekt, den kommer fråga dig om projektets namn osv. När du är klar med det så handlar det om att lägga in såkallade `Targets' i projectet. Ett Target är en målfil som skall skapas, som tex databas.o eller liknande. För att skapa ett nytt Target skriver du M-x ede-new-target RET Sedan mata in vad ditt target ska heta. Vanligtvis frågar programmet om du vill lägga till filen du står i men du kan även lägga till nya filer senare med kommandot M-x ede-add-file för att lägga till fler beroenden till din målfil.

När du har ett projekt, och minst ett Target i ditt projekt kan du bygga projektet genom snabbkommandot C-c . C då generar programmet automatiskt en Makefil och kör den , det finns öven stöd för att kompilera delmål och generera om Makefilen  

\newpage
\section{ECB (Emacs Code Browser)}

\subsection{Installation och Konfigurering av ECB}
ECB behöver Cedet för att fungera korrekt, så du måste installera det först om du inte har det installerat redan, vissa nyare versioner har cedet redan installerat. Men för att installera ECB på lättaste sättet så öppna upp din ~/.emacs eller ~/.emacs.d/init.el och lägg in följande rader i filen
\begin{lstlisting}
(setq package-archives 
'(("gnu" . "http://elpa.gnu.org/packages/")
  ("marmalade" . "http://marmalade-repo.org/packages/")
  ("melpa" . "http://melpa.milkbox.net/packages/")))
\end{lstlisting}
spara filen och starta om emacs, evt skriv in M-x eval-buffer RET. Detta lägger till att arkiv av paket, som sedan går att installera.

Nästa steg är att skriva M-x list-packages RET. Då kommer det en lång lista med olika paket. Sök efter ECB antingen genom att scolla listan eller tryck C-s ecb , så borde du hitta raden. Tryck på ecb ordet och det bör öppnas en ny buffer. I den nya buffertern, tryck på Install och låt scriptet köras.


OBS!: Härmed ska programmet vara installerat, men du kan ha fått upp varingsmeddelanden. Detta kan bero på olika saker men först så navigera dig till ~/.emacs.d/elpa/ecb-20130826.1941/ecb-cedet-wrapper.el eller motsvarande sökväg. I filen navigera till raden där det står \begin{lstlisting}
(defconst ecb-cedet-required-version-max '(x x x x)
\end{lstlisting}
och byt ut mot följande rad
\begin{lstlisting}
(defconst ecb-cedet-required-version-max '(2 1 4 9)
\end{lstlisting}
spara sedan filen och starta om emacs.

När du startar upp emacs igen så ska det fungera genom att skriva M-x ecb-activate RET eller M-x ecb-minor-mode RET. Du kan med fördel lägga till raden med paranterser om sig i din .emacs fil om du vill.
\newpage
\subsection{Fönster i ECB}
ECB består av en gäng olika fönster som fyller olika funktioner för användaren som underlättar navigering i och utanför din kod.
\begin{description}
\item [ECB Directories-window (C-c . g d)] \hfill \\
  Visar en trädstruktur från din satta source path.
\item [ECB Source-window (C-c . g s)] \hfill \\
  Visar filer i katalogen valda katalogen i ditt Directory-Window
\item [ECB Method-window (C-c . g m)] \hfill \\
  Visar en överblick på ditt aktiva document och ger dig möjlighet att hoppa till funktioner och variabler osv. i din kod. (om du har stöd för att parseing av språket, tex för Java behövs JDEE installerat)  
\item [ECB History-window (C-c . g h)] \hfill \\
  Visar vilka filer du har haft öppna och ger dig möjlighet att hoppa mellan på ett smidigt sätt.
\item [ECB Edit-Window 1 (C-c . g e 1)] \hfill \\
  Visar din aktuella buffer som du arbetar i.(Om du jobbar med flera Edit fönster uppe kan du byta ut siffran 1:an till 2,3 ... )
\item [ECB Compile-Window (C-C . g c)] \hfill \\
  Visar din compile buffer, när du gör kommandon som compile. C-c . / ger dig lite större buffer att läsa ur och C-c . \ öppnar och stänger fönstret
\end{description}

\subsection{Navigering med ECB}
Navigering i görs med antingen tangentbordet med komandona som är beskrivna ovan i kombination med piltangenterna och return-knappen eller med musen. Det som kan vara bra att notera är att det finns något som kallar för Power Click som du kan med hjälp av shift knappen och sedan antingen return eller musklick så uppdateras din katalog om den av någon anledning har blivit utdaterad. Och för att expandera en nod i någon av trädstrukturerna tycker du bara på TAB knappen så öppnas eller stängs fliken. Det finns mängder med kommandon du kan köra som finns beskrivna på hemsidan \url{http://ecb.sourceforge.net} där du läsa mer om att konfigura din miljö precis som du vill ha den.
\subsection{Ändra layout i ECB}
Det finns fördefinerade layouter för ECB som man kan använda sig av om man vill ha en annan struktur eller kanske bara vissa fönster öppna. Du kommer lättast åt det genom att skriva M-x ecb-change-layout RET , då frågar programmet dig om ett namn på layouten du vill ha då kan du lätt bara tycka TAB så dyker en buffer upp med de olika fördefinerade möjligheterna. Testa dig runt och se om det är någon som passar dig bra.

%----------------------------------------------------------------------------------------
%	MAJOR SECTION X - TEMPLATE - UNCOMMENT AND FILL IN
%----------------------------------------------------------------------------------------

%\section{Content Section}

%\subsection{Subsection 1} % Sub-section

% Content

%------------------------------------------------

%\subsection{Subsection 2} % Sub-section

% Content

%----------------------------------------------------------------------------------------
%	CONCLUSION
%----------------------------------------------------------------------------------------
\section{Sammanfattning}
Självständigt är Emacs ett kraftfullt verktyg som kan det mesta, men med dessa verktyg kommer man steget längre att smidigt navigera sig i koden och ger möjligheter att överskåda och arbeta tillsammans på projekt. 
Det finns många funktioner att undersöka i både Cedet och ECB som kan göra ditt arbete i emacs blir ännu mer produktivt som tillexempel hur man generar distubutioner av din kod eller laddar in ett projekt som någon annan byggt och ger dig möjligheten att sammarbeta i samma projekt. Finns även stöd olika typer av versions kontroll som kan göra ett sammarbetesprojekt extremt mycket smidigare 

\section{Lära Mer} 
Det finns alltid rum för att lära sig mer och utforka verktygen till det yttersta. Detta har varit en introduktion till användandet av Cedet och ECB men du kan alltid finna mer information och dokumentation på deras hemsidor
\begin{itemize}
\item{CEDET} \url{http://cedet.sourceforge.net}
\item{ECB} \url{http://ecb.sourceforge.net}
\end{itemize}


%----------------------------------------------------------------------------------------

\end{document}
